\section{Rational lines and divisorial decompositions}\label{sec:preliminaries}

The proof keeps two kinds of information separate: a real \((1,1)\)-class
and the rational holomorphic line bundle representing it. We first set
out this distinction and the precise ordinary and projective results used
below.

\subsection{Conventions and changes of model}

We work over \(\C\). A rational line bundle on a complex space \(X\) means
an element of \(\Pic(X)\otimes_{\Z}\Q\); we use additive notation.
The relation \(L\qsim L'\) is equality in this group, so that sufficiently
divisible integral multiples are isomorphic holomorphic line bundles.
On a smooth compact K\"ahler manifold, \(L\num L'\) means
\(c_1(L)=c_1(L')\) in real Bott--Chern cohomology. We also write
\(\{E\}=c_1(\cO_X(E))\) for a real divisor \(E\).
The \(\partial\bar\partial\)-lemma identifies real Bott--Chern
\((1,1)\)-classes with real de Rham classes of type \((1,1)\).

A class is \emph{pseudo-effective} if it contains a closed positive current;
it is \emph{nef} if it lies in the closure of the K\"ahler cone.
A rational line is \emph{semiample} if a positive integral multiple is
generated by its global sections. On a smooth projective variety these
notions, for divisor classes, agree with the corresponding algebraic
ones, and equality of real divisor classes agrees with numerical
equivalence. We use the analytic definitions on every nonprojective
model.

All manifolds and normal spaces are connected unless otherwise stated.
A modification is a proper bimeromorphic morphism. Resolutions,
flattenings, and main components of fiber products are always followed,
when necessary, by normalization and a smooth compact K\"ahler
modification. Spaces in Fujiki class \(\mathcal C\) are used only through
such models. Projective targets admit projective modifications. The
resolution procedures can be chosen projective and functorial; finite
covers of compact K\"ahler spaces are K\"ahler.

We use the usual discrepancy definition of a klt pair, with an effective
rational boundary unless a signed boundary is expressly mentioned.
A signed pair with all discrepancies greater than \(-1\) is called
sub-klt. On a smooth model, a simple normal crossing boundary is klt
precisely when all its coefficients are less than one.

For a normal klt pair \((X,B)\) and a log resolution \(\mu:\widetilde X\to X\),
write
\[
 K_{\widetilde X}+\widetilde B=\mu^*(K_X+B).
\]
Replacing \(\widetilde B\) by its coefficientwise positive part produces
an effective simple normal crossing klt boundary \(B^+\), and
\begin{equation}\label{pre:resolution-error}
 K_{\widetilde X}+B^+=\mu^*(K_X+B)+E,
 \qquad E\geq0\text{ exceptional over }X.
\end{equation}
The exceptional support assertion uses effectiveness of \(B\): the
nonexceptional coefficients of \(\widetilde B\) are already nonnegative.
The pullback of a nef rational line remains nef. Thus this convention
preserves the hypotheses of the induction.

\begin{lemma}[Lines of zero class]\label{pre:picard}
Let \(\mu:Y\to X\) be a modification between smooth compact K\"ahler
manifolds.
\begin{enumerate}
\item A rational line whose real Chern class is zero belongs to
\(\Pic^0(X)\otimes_{\Z}\Q\).
\item Pullback identifies \(\Pic^0(X)\) with \(\Pic^0(Y)\).
\item If a rational line \(L\) on \(X\) has semiample pullback, then \(L\)
is semiample.
\end{enumerate}
\end{lemma}
\begin{proof}
After clearing denominators, a line of zero real class has torsion
integral Chern class. Another positive power has zero integral class,
so belongs to \(\Pic^0\) by the exponential sequence.
The second assertion is the standard bimeromorphic invariance of
\(\Pic^0\): a smooth modification preserves \(H^1(\cO)\) and the
integral lattice defining this torus. It also follows by factoring
smooth bimeromorphic maps into blowups and blowdowns with smooth centers.
Finally \(\mu_*\cO_Y=\cO_X\) by normality. The projection formula identifies
the sections of every integral multiple of \(L\) with those of its
pullback. A base point downstairs would therefore be a base point at
every point above it.
\end{proof}

\subsection{The analytic negative part}

For a pseudo-effective real class \(\alpha\) on a smooth compact
K\"ahler manifold, fix a K\"ahler form \(\omega\). If \(E\) is a prime
divisor, its minimal multiplicity is
\[
 \nu_E(\alpha)=\lim_{\varepsilon\downarrow0}
 \inf\{\nu_E(T):T\in\alpha,\quad T\geq-\varepsilon\omega\}.
\]
Here \(\nu_E(T)\) is the generic Lelong number. The definition is
independent of \(\omega\), and currents with analytic singularities
suffice. Boucksom's divisorial decomposition is
\[
 N(\alpha)=\sum_E\nu_E(\alpha)E,\qquad
 Z(\alpha)=\alpha-\{N(\alpha)\}.
\]
Its negative part is an effective real divisor, and its positive part
is \emph{modified nef}: all its minimal divisorial multiplicities vanish.
For a rational line \(L\), write \(N(L)=N(c_1(L))\).
The foundational analytic results are due to Boucksom
\cite[Sections~2--3 and~5]{Boucksom2004}. The following forms, including
their use on smooth resolutions of normal K\"ahler spaces, are recorded
and proved in \cite[Section~2]{K}.

\begin{lemma}[Negative-part calculus]\label{pre:negative}
Let \(\alpha\) be pseudo-effective on a smooth compact K\"ahler manifold.
\begin{enumerate}
\item Every positive current in \(\alpha\) contains
\([N(\alpha)]\). If \(0\leq F\leq N(\alpha)\), then
\[
 N(\alpha-\{F\})=N(\alpha)-F.
\]
\item If \(\beta\) is nef, then \(N(\alpha+\beta)\leq N(\alpha)\).
In particular \(N(\beta)=0\).
\item A modified nef class restricts pseudo-effectively to a resolution
of every prime divisor.
\item If \(\mu:Y\to X\) is a smooth modification and
\(\alpha=\beta+\{R\}\), with \(\beta\) nef and \(R=N(\alpha)\), then
\[
 N(\mu^*\alpha)=\mu^*R.
\]
\item Suppose \(\mu:Y\to X\) is a modification from a smooth compact
K\"ahler manifold to a normal compact K\"ahler space, \(A\) is a rational
line on \(X\), and \(E\geq0\) is \(\mu\)-exceptional. If
\(\mu^*A+E\) is pseudo-effective, then \(\mu^*A\) is pseudo-effective and
\[
 N(\mu^*A+E)=N(\mu^*A)+E.
\]
\end{enumerate}
\end{lemma}

These are \cite[Lemmas~2.2--2.4]{K}. The exceptional translation in
part~(5) also holds for a real class with smooth local potentials on
the normal target. Part~(3) concerns prime divisors; we do not use
unrestricted pseudo-effective restriction to arbitrary subvarieties.
Section~\ref{sec:completion} gives the related intersection argument
needed when a map has positive-dimensional fibers.

Two consequences will be used repeatedly. An effective divisor
representing a multiple \(mL\) contains \(mN(L)\), because its divisorial
current is positive. Also, the decomposition sought in
Theorem~\ref{thm:decomposition} is unchanged by taking a higher smooth
model: part~(4) pulls it up, and part~(5) removes the error
\eqref{pre:resolution-error}.

\subsection{The ordinary K\"ahler input}

The first substantial input is ordinary log abundance, in its
divisorial form. We state the full line-bundle information needed
in the proof. The restricted model and contraction package of
\cite[Theorem~3.55]{K} uses the base presentations of
\cite[Proposition~3.30]{K} for the whole descended adjoint class.
At its generalized klt (gklt) contraction and relative-model steps,
these presentations give gklt base data with a globally nef carrier
datum and modified-big
total boundary. At the non-klt contraction step, the multiplier-ideal
inclusion, projective descent, and restriction supply the required
contraction on the full non-klt closed subspace. The ordinary adjoint
induction then yields the following actual rational-line statement.

\begin{theorem}[Ordinary divisorial decomposition]\label{pre:ordinary}
Let \(X\) be a smooth compact K\"ahler manifold, and let \(B\) be a
rational simple normal crossing divisor with coefficients in \([0,1]\).
If \(J=K_X+B\) is pseudo-effective, there is a smooth compact K\"ahler
modification \(\mu:U\to X\) and an identity
\[
 \mu^*J\qsim P+R,\qquad P\text{ semiample},\qquad
 R=N(\mu^*J)\geq0,
\]
where \(P\) is a rational line and \(R\) is a rational divisor.
\end{theorem}

This is \cite[Theorem~2.13]{K}, whose logarithmic Iitaka hypothesis is
supplied by \cite[Corollary~6.2]{LI}. More precisely, the required inequality
is
\[
 \kappa(X,K_X+D_X)\geq
 \kappa(F,K_F+D_X|_F)+\kappa(Y,K_Y+D_Y)
\]
for a surjective morphism \(f:X\to Y\) with connected fibers between
smooth connected projective varieties, a very general smooth fiber
\(F\), reduced simple normal crossing boundaries \(D_X,D_Y\),
allowing zero, and \(\Supp(f^*D_Y)\subseteq\Supp(D_X)\).
The conventions are \((-\infty)+b=-\infty\) and Iitaka dimension
zero for a point. This is precisely \cite[Assumption~1.1]{K}, so
\cite[Corollary~6.2]{LI} supplies the hypothesis in every finite
dimension. We use Theorem~\ref{pre:ordinary} in every dimension. It implies ordinary nonvanishing: a sufficiently divisible
multiple of \(J\) has a nonzero section. If \(a(X)=0\), its semiample
part is rationally trivial, and hence \(\mu^*J\qsim R\).

We also use the following program consequence of ordinary decomposition.
The version stated here is confined to its smooth starting models.

\begin{proposition}[Contraction of a known negative part]\label{pre:contract}
Let \((X,B)\) be a smooth compact K\"ahler klt pair with rational
simple normal crossing boundary. Suppose
\[
 K_X+B\qsim P+R,\qquad P\text{ semiample},\qquad
 R=N(K_X+B)\geq0
\]
with \(P,R\) rational. A finite ordinary \((K_X+B)\)-negative program
reaches a normal compact K\"ahler klt model on which the log canonical
line is semiample. Its steps contract or flip detected analytic extremal
rays, are \(P\)-trivial, and preserve a fixed generated Cartier multiple
of \(P\) as an actual line. The final adjoint is rationally linearly
equivalent to the descended semiample line. The steps extract no
divisors; on a common resolution the initial adjoint equals the
pullback of the final adjoint plus an effective divisor exceptional
over the final model.
\end{proposition}

This is the ordinary klt case of
\cite[Proposition~3.8]{K}, using the ordinary scaling construction of
\cite[Lemma~3.6]{K} and its detected analytic rays. Smoothness supplies
global strong \(\Q\)-factoriality, the displayed identity and
Lemma~\ref{pre:negative}(4) supply its resolution hypothesis, and
Theorem~\ref{pre:ordinary} supplies the lower-dimensional ordinary
hypotheses in that proposition.
The ray estimate used to keep this program over a prescribed
projective base will be verified at the point of use.

\subsection{The projective input}

The second substantial input is numerical semiampleness of the positive
part for projective klt pairs. It is important to apply it to a nef
line on the original model; the trace of its b-divisor on a later
minimal model need not remain nef.

\begin{proposition}[Projective positive parts]\label{pre:projective}
Let \((S,\Delta)\) be a normal projective klt pair with rational boundary,
and let \(M\) be a nef rational Cartier divisor. Suppose
\(K_S+\Delta\) is pseudo-effective.
There is a smooth projective modification \(u:S'\to S\), a birational
morphism \(v:S'\to S_m\) to a normal projective variety, a nef rational
Cartier divisor \(H_m\) on \(S_m\), and an effective rational
\(v\)-exceptional divisor \(E\), such that
\begin{equation}\label{pre:projective-model}
 u^*(K_S+\Delta+M)\qsim v^*H_m+E.
\end{equation}
On a sufficiently high such model, \(v^*H_m\) is numerically equivalent
to a semiample rational divisor and
\[
 N\bigl(u^*(K_S+\Delta+M)\bigr)=E.
\]
\end{proposition}
\begin{proof}
The positive-part assertion of \cite[Proposition~8.1]{P} gives a smooth
model on which the rational positive part of the original adjoint
is numerically semiample. Separately, take a small projective
\(\Q\)-factorialization and run a terminating generalized klt program
as in \cite[Theorem~2.3]{P}. Its fixed nef data are the pullback of \(M\).
The generalized discrepancies initially agree with the ordinary ones
because \(M\) descends, and the generalized adjoint is pseudo-effective.
The resulting model has nef adjoint \(H_m\). The comparison on a common
resolution is \eqref{pre:projective-model}, with effective exceptional
error.

Choose this resolution also to dominate the model furnished by
the positive-part theorem. Lemma~\ref{pre:negative}(5) and nefness of
\(v^*H_m\) identify \(E\) as the negative part. Pulling up the first
model's decomposition by Lemma~\ref{pre:negative}(4) identifies its
numerically semiample positive class with \(v^*H_m\). Algebraic and
analytic negative parts agree on these smooth projective models.
\end{proof}

The generalized program used in this proof is the fixed-nef-data
minimal-model theorem cited in \cite{P}; it is not an application of
the numerical semiampleness assertion being proved here. The two inputs
above will be used as stated theorems. The remaining sections prove the
additional K\"ahler arguments, including the descent and period
constructions that permit their application.
